% =====================================================================
% Sekcja 2 - Tło teoretyczne (wersja kompaktowa, ~1 strona)
% =====================================================================

\section{Tło teoretyczne}
\label{sec:tlo}

\subsection{Perplexity jako miara modelu językowego}
\label{sec:perplexity-def}

Perplexity (PPL) to średnia geometryczna odwrotności prawdopodobieństw,
jakie model przypisuje kolejnym tokenom danej sekwencji. Dla
sekwencji $x_1, \dots, x_N$ pod modelem $p_\theta$:
\begin{equation}
\mathrm{PPL}(x_{1:N}) = \exp\!\left(-\frac{1}{N} \sum_{t=1}^{N} \log p_\theta(x_t \mid x_{<t})\right).
\label{eq:ppl}
\end{equation}
Niższa wartość oznacza, że model był mniej zaskoczony tym, co
zobaczył. Klasycznie PPL używa się jako miary jakości samego
modelu językowego - im niższe na zbiorze testowym, tym lepiej
model przewiduje tekst. W tej pracy interesuje jednak inne
zastosowanie: PPL liczone na własnej odpowiedzi modelu, jako
sygnał mówiący o tym, czy ta odpowiedź jest poprawna - taki,
żeby agent mógł np.\ odroczyć odpowiedź lub poprosić o pomoc,
gdy jego pewność spada.

\subsection{Argumenty za perplexity}
\label{sec:pro-ppl}

Gonen i in.~\cite{gonen2023demystifying} pokazali, że perplexity
samego promptu jest dobrym predyktorem skuteczności modelu - prompty
o niższym PPL dają wyższą trafność odpowiedzi. Interpretacja
autorów: prompty, które ,,brzmią naturalniej'' dla modelu (mniejsze
zaskoczenie na jego tokenach), są dla niego łatwiejsze i prowadzą
do lepszych wyników. Sam pomiar PPL działa tu jako tania heurystyka
wyboru lepszego promptu - bez konieczności faktycznego uruchamiania
zadania i sprawdzania wyniku.

\subsection{Argumenty przeciw perplexity}
\label{sec:anti-ppl}

Wang i in.~\cite{wang2022perplexity} pokazują drugą stronę medalu:
PPL bywa zawodne jako miara jakości tekstu. Autorzy zwracają uwagę
na trzy pułapki. Po pierwsze, krótsze teksty dostają wyższe PPL
niezależnie od jakości (\emph{length bias}). Po drugie, powtarzające
się fragmenty sztucznie je zaniżają (\emph{repetition sensitivity}).
Po trzecie, drobne różnice w interpunkcji potrafią wprowadzać dużą
zmienność (\emph{punctuation dependency}). Zestawienie tych dwóch
prac otwiera pytanie istotne dla tego raportu: czy ten obraz
zmienia się, gdy zamiast pojedynczego promptu analizowane są
wieloetapowe trajektorie agentów?

\subsection{Perplexity w długim kontekście}
\label{sec:longppl}

Fang i in.~\cite{fang2025longppl} pokazują problem klasycznego PPL
na zadaniach z długim kontekstem. Standardowe PPL uśrednia
logarytmy prawdopodobieństw po wszystkich tokenach z taką samą
wagą, przez co całość zdominowana jest przez tokeny lokalnie
przewidywalne - spójniki, końcówki fleksyjne, znaki interpunkcyjne -
a wkład tokenów, które naprawdę zależą od dalekiego kontekstu,
ginie w średniej. Zaproponowana przez autorów metryka LongPPL
waży tokeny zgodnie z tym, jak bardzo zależą od dalekiego kontekstu;
na benchmarkach długiego kontekstu osiąga ona korelację Pearsona
$r \approx -0{,}96$ z faktyczną wydajnością modeli - podczas gdy
zwykłe PPL daje na tych samych danych wartość bliską zeru.

Ten sam mechanizm może działać też w przypadku trajektorii agentów.
Trajektorie agenta ReAct - złożone z naprzemiennych bloków
\texttt{Thought}, \texttt{Action}, \texttt{Observation}
i \texttt{Answer} - są długie i pełne tokenów strukturalnych
(samych etykiet segmentów, znaków interpunkcyjnych), na których
model jest bardzo pewny niezależnie od jakości rozwiązania. Takie
tokeny mogą zdominować jedną wartość PPL liczoną na całej
trajektorii i wytłumić sygnał z miejsc, w których model
rzeczywiście ,,myśli'' - na przykład formułuje wywołanie
kalkulatora. Stąd w tej pracy oprócz PPL na całej trajektorii
liczone jest też PPL osobno na segmentach \texttt{Action}
i \texttt{Answer} (sprawdzane w sekcji~\ref{sec:wyniki-agent}).
Szerszy przegląd metryk ewaluacji agentów - w którym PPL pojawia
się obok trafności wywołań narzędzi, jakości planu, metryk
trajektorii i bezpieczeństwa - można znaleźć w pracy Mohammadi
i in.~\cite{mohammadi2025survey}.

\subsection{Niepewność jako sygnał sterujący w agentach}
\label{sec:agent-uncertainty}

Trzy nowsze artykuły idą krok dalej: zamiast tylko mierzyć pewność
modelu, wpinają ją wprost w pętlę decyzyjną agenta. ReDAct~\cite{redact2026}
wprowadza próg pewności, poniżej którego agent odracza odpowiedź
zamiast ryzykować błąd. \emph{Every Response Counts}~\cite{everyresponse2026}
zbiera niepewność wielu agentów pracujących nad tym samym zadaniem
i - korzystając z dekompozycji tensorowej - próbuje wyłowić, kiedy
modele systematycznie się ze sobą nie zgadzają. \emph{Towards
Agents That Know When They Don't Know}~\cite{towardsknow2025}
traktuje niską pewność jako sygnał, że trzeba się zatrzymać
i sprawdzić aktualny krok rozumowania, zanim agent pójdzie dalej.

Wszystkie trzy prace milcząco zakładają, że pewność liczona
z rozkładu tokenów rzeczywiście mówi coś o jakości tego, co model
generuje. W tej pracy sprawdzane jest, na ile to założenie się
utrzymuje na prostym agencie ReAct rozwiązującym zadania matematyczne:
jak silna jest korelacja PPL i entropii z poprawnością i czy ma
znaczenie, na której części trajektorii się ją mierzy.
